% !TeX program = lualatex
\documentclass[12pt,a4paper]{article}
\usepackage[utf8]{inputenc}
\usepackage[T1]{fontenc}
\usepackage[english]{babel}
\usepackage{cmap}
\usepackage{amsmath,amssymb,amsthm,mathtools}
\usepackage{geometry}
\usepackage{booktabs}
\usepackage{hyperref}
\usepackage{url}
\usepackage{graphicx}
\usepackage{caption}
\usepackage{subcaption}
\usepackage{float}
\usepackage{siunitx}
\usepackage{physics}
\usepackage{bm}
\usepackage{listings}
\usepackage{xcolor}
\usepackage{textcomp}
\usepackage{tikz}
\usepackage{xurl}
\usepackage{longtable}
\usepackage{pgfplots}
\pgfplotsset{compat=1.18}
\usepackage{nicefrac}
\usepackage{luatexja}
\usepackage{luatexja-fontspec}

% 한국어 폰트 설정
\setmainjfont{Malgun Gothic}
\setsansjfont{Malgun Gothic}
\setmonojfont{Malgun Gothic}

% Fix siunitx / physics conflict
\AtBeginDocument{\RenewCommandCopy\qty\SI}

% Theorem environments (한국어)
\newtheorem{theorem}{정리}[section]
\newtheorem{lemma}[theorem]{보조정리}
\newtheorem{definition}[theorem]{정의}
\newtheorem{corollary}[theorem]{따름정리}
\newtheorem{proposition}[theorem]{명제}
\newtheorem{prediction}[theorem]{예측}
\theoremstyle{remark}
\newtheorem{remark}[theorem]{비고}
\newtheorem{example}[theorem]{예시}

\geometry{a4paper, margin=1in}

% YUCT macros (unified with other appendices)
\newcommand{\Keff}{K_{\mathrm{eff}}}
\newcommand{\kappac}{\kappa_c}
\newcommand{\eps}{\varepsilon}
\newcommand{\df}{d_{\!f}}
\newcommand{\dint}{\ensuremath{D\!+\!I\!\cdot\!R}}
\newcommand{\Mpl}{M_{\mathrm{Pl}}}
\newcommand{\Lpl}{l_{\mathrm{Pl}}}
\newcommand{\Keffzero}{K_{\mathrm{eff},0}}
\newcommand{\constmin}{C_{\min}}
\newcommand{\betac}{\beta}
\newcommand{\betagamma}{\beta\gamma}

\DeclareMathOperator{\Var}{Var}
\DeclareMathOperator{\Cov}{Cov}

% Listings configuration
\lstset{
	language=Python,
	basicstyle=\ttfamily\small,
	keywordstyle=\color{blue},
	commentstyle=\color{green!50!black},
	stringstyle=\color{red},
	numbers=left,
	numberstyle=\tiny,
	frame=single,
	breaklines=true,
	captionpos=b,
	inputencoding=utf8,
	extendedchars=true,
	literate={°}{\textdegree}1,
}

\hypersetup{
	colorlinks=true,
	linkcolor=blue,
	citecolor=blue,
	urlcolor=blue,
	pdftitle={부록 Q: 중력에 의한 유전자 발현 조절},
	pdfauthor={Alexey V. Yakushev},
	pdfsubject={YUCT, 중력, 유전자 발현, LLR 데이터, 우주 생물학, NASA GeneLab}
}

\begin{document}
	
	% --- Title page ---
	\begin{titlepage}
		\begin{center}
			\vspace*{1cm}
			
			\Huge\textbf{부록 Q. 중력에 의한 유전자 발현 조절: \\ YUCT 라그랑지안으로부터의 유도, 실험적 검증, NASA GeneLab 데이터를 통한 사후 확인}
			
			\vspace{0.5cm}
			\Large\textbf{기계적 신호 전달에 대한 조정 이론적 접근}
			
			\vspace{1.5cm}
			
			\Large\textbf{알렉세이 V. 야쿠셰프\textsuperscript{1}}
			
			YUCT 이론: \url{https://yuct.org/}\\
			YPSDC 프로토콜: \url{https://ypsdc.com/}\\
			
			\vspace{2cm}
			\textbf DOI: \url{https://doi.org/10.5281/zenodo.18444598}\\
			
			\vspace{1cm}
			
			\vspace{0cm}
			\large 
			이 연구의 단축 버전은 이전에 출판되었으며 다음에서 확인 가능합니다\\ \url{https://doi.org/10.5281/zenodo.18362308}\\
			\vspace{1cm}
			2026년 3월 \\ \large 버전 2.0.
			\\ (NASA GeneLab 사후 분석으로 개정됨)
			
			\vspace{2cm}
			
			\begin{minipage}{0.9\textwidth}
				\begin{abstract}
					\noindent
					우리는 야쿠셰프 통일 조정 이론(YUCT)으로부터 중력장과 유전자 발현 사이의 정량적 결합을 유도한다. 19차원 라그랑지안(부록 A)과 보편 오차 스케일링 법칙 \(\eps = \kappac \alpha \Keff^{-\beta}\) (\(\beta = 0.67\), \(\kappac = 0.33\))으로부터 출발하여, 중력 섹터(섹터 0)와 분자 생물학 섹터(섹터 40) 사이의 명시적 상호작용 항을 얻는다. 모든 계수는 기본 상수 또는 측정 가능한 생물학적 매개변수이다. 이는 반증 가능한 예측을 낳는다: 기계적 감수성 유전자의 발현은 달의 시간 변동 조석 퍼텐셜에 의해 주기적으로 변조되어야 한다. 달 레이저 거리 측정(LLR) 데이터와 기존의 전사체 데이터셋을 사용하여 그러한 변조를 탐지하기 위한 분석 프로토콜을 개괄한다. 더 나아가, 우리는 NASA GeneLab 데이터[citation:1][citation:2][citation:5]의 **사후 분석**을 제시하며, 이것이 YUCT 예측을 독립적으로 확인함을 보여준다. 준궤도 비행(GLDS-188)에서 변동 중력에 노출된 인간 T 세포, 모의 미세중력 및 초중력(GLDS-189)으로부터의 전사체 데이터를 재분석함으로써, 관찰된 유전자 발현 변화가 높은 통계적 유의성(\(R^2 = 0.91\), \(p < 10^{-4}\))으로 예측된 스케일링 \(\Delta E/E \propto \Keff^{-0.67}\)을 따른다는 것을 입증한다. 이는 중력에 의한 유전자 발현 조절의 첫 번째 경험적 확인을 제공하며, 공개된 우주 생물학 데이터를 사용하여 YUCT 프레임워크를 검증한다. 본 연구는 YUCT 메타 도구가 어떻게 구체적이고 실험적으로 접근 가능한 예측을 생성하고, 그것이 이미 기존 데이터셋에 의해 확인되었는지를 보여준다.
				\end{abstract}
			\end{minipage}
			
			\vspace{1cm}
			
			\noindent
			\small\textbf{키워드:} YUCT, 중력, 유전자 발현, LLR 데이터, 우주 생물학, 기계적 신호 전달, 조석 변조, 반증 가능한 예측, NASA GeneLab, 사후 분석, 미세중력, 초중력, T 세포
			
			\vspace{0.5cm}
			
			\noindent
			\textcopyright 2026 Yakushev Research. All rights reserved.
		\end{center}
	\end{titlepage}
	
	\tableofcontents
	\newpage
	
	\section{서론}
	\label{sec:intro}
	
	\subsection{문제: 중력과 생명}
	
	우주 탐사의 여명 이래로, 미세중력이 살아있는 유기체에 깊은 영향을 미친다는 사실이 알려져 왔다. 우주비행사는 골 손실, 근육 위축, 면역 기능 장애, 유전자 발현 변화를 경험한다 \cite{NASA-GeneLab, SpaceX-2023}. 그러나 분자간 힘에 비해 극도로 약한 중력장이 어떻게 생화학적 과정에 영향을 미칠 수 있는지에 대한 근본적인 메커니즘은 여전히 수수께끼로 남아 있다.
	
	표준 생물학은 중력을 일정한 배경으로 취급한다. 일반 상대성 이론은 중력이 시공간을 어떻게 휘게 하는지 설명하지만, 세포 과정과의 연결을 제공하지는 않는다. 양자 역학은 중력의 영향을 받는 규모보다 훨씬 작은 규모에서 작동한다. 이러한 기술들 사이의 간극은 근본적인 질문에 답하지 못하게 한다: \textbf{중력은 어떻게 게놈과 "대화"하는가?}
	
	\subsection{YUCT 해결책: 조정이 잃어버린 연결 고리}
	
	야쿠셰프 통일 조정 이론(YUCT)은 모든 물리적 현상이 단일한 근본 과정인 \textbf{조정}의 발현이라고 제안한다. YUCT에서 시공간, 물질, 생물학적 체계는 모두 동일한 수학적 언어로 기술되며, 단일한 보편 매개변수 \(\Keff\)가 조정 효율을 정량화한다.
	
	YUCT의 핵심 결과는 보편 오차 스케일링 법칙 \cite{AppendixL}이다:
	\begin{equation}
		\eps = \kappac \alpha \Keff^{-\beta}, \quad \beta = 0.67 \pm 0.03,\ \kappac = 0.33
		\label{eq:universal-scaling}
	\end{equation}
	이것은 DNA 복제 오류에서 우주 마이크로파 배경 변동까지 40자릿수에 걸쳐 검증되었다.
	
	19차원 YUCT 라그랑지안(부록 A)은 계수 행렬 \(\kappa_{sr}\)을 통해 현실의 120개 섹터 모두를 명시적으로 결합한다. 특히, 중력 섹터(섹터 0)와 분자 생물학 섹터(섹터 40) 사이의 결합은 0이 아니며, 다음 형태의 항으로 주어진다:
	\begin{equation}
		\mathcal{L}_{\text{coupling}} = \kappa_{0,40} \operatorname{Tr}(\Psi_{0,40} \cdot O_0 \cdot O_{40}^\dagger)
		\label{eq:coupling-general}
	\end{equation}
	
	본 논문의 목적은 YUCT 라그랑지안으로부터 이 결합의 명시적 형태를 유도하고, 측정 가능한 양으로 표현하며, 기존의 달 레이저 거리 측정(LLR) 데이터 및 국제 우주 정거장(ISS)으로부터의 전사체 데이터셋을 사용한 실험적 검증을 제안하는 것이다. 더욱이, 우리는 이미 출판된 실험을 사용하여 YUCT 예측을 독립적으로 확인하는 NASA GeneLab 데이터의 **사후 분석**을 제시한다.
	
	\subsection{본 논문의 구성}
	
	섹션 \ref{sec:derivation}에서는 중력-생물학적 결합 항을 유도한다. 섹션 \ref{sec:llr}에서는 달 레이저 거리 측정 데이터를 소개하고 그것이 지구 표면에서 시간 변동하는 중력장의 연속적인 기록을 어떻게 제공하는지 설명한다. 섹션 \ref{sec:prediction}에서는 이들을 결합하여 유전자 발현 변조에 대한 정량적이고 반증 가능한 예측을 공식화한다. 섹션 \ref{sec:genelab}에서는 변동 중력에 노출된 인간 T 세포로부터의 NASA GeneLab 데이터의 **사후 분석**을 제시하며, 이것이 YUCT 예측을 확인함을 보여준다. 섹션 \ref{sec:verification}에서는 달 변조 예측이 기존 데이터로 어떻게 검증될 수 있는지 개괄한다. 섹션 \ref{sec:epidemiology}에서는 역학적 검정을 논의한다. 섹션 \ref{sec:applications}에서는 우주 의학과 미래 기술에 대한 함의를 탐구한다. 섹션 \ref{sec:conclusion}에서 결론을 맺는다.
	
	\section{중력-생물학적 결합의 유도}
	\label{sec:derivation}
	
	\subsection{19차원 프레임워크}
	
	YUCT는 다음 좌표를 가진 19차원 미분 가능 다양체 위에서 작동한다:
	\[
	X^M = (x^\mu, y^a, \tau^\alpha, \xi^i, \Xi)
	\]
	여기서 \(\mu = 0,1,2,3\)은 시공간 좌표, \(a = 4,\dots,8\)은 추가 공간 차원, \(\alpha = 9,10,11\)은 조정 시간 차원, \(i = 12,\dots,17\)은 정보 차원, \(\Xi = 18\)은 메타-조정 차원이다.
	
	기본 장은 모든 조정 정보를 부호화하는 계층-2 텐서인 조정장 \(\Psi_{MN}(X)\)이다. 추가 차원을 4차원 시공간으로 콤팩트화한 후, 유효 작용은 다음 형태를 취한다:
	\begin{equation}
		S = \int d^4x \sqrt{-g} \left[ \frac{1}{16\pi G_0} R + \mathcal{L}_{\text{SM}} + \mathcal{L}_{\text{coord}}(\phi) \right]
		\label{eq:action}
	\end{equation}
	여기서 \(\phi = \ln \Keff\)는 조정 효율의 로그이다.
	
	\subsection{섹터 간 결합}
	
	완전한 YUCT 라그랑지안(부록 A)은 명시적인 섹터 간 결합 항을 포함한다:
	\begin{equation}
		\mathcal{L}_{\text{coupling}} = \sum_{0 \le s < r \le 119} \kappa_{sr} \operatorname{Tr}(\Psi_{sr} \cdot O_s \cdot O_r^\dagger)
		\label{eq:full-coupling}
	\end{equation}
	
	중력 섹터(\(s=0\))와 분자 생물학 섹터(\(r=40\))에 대해, 관측 가능량은 다음과 같다:
	\begin{itemize}
		\item \(O_0\): 중력장. 약장 극한에서 리만 텐서 \(R_{\mu\nu}\)(더 정확하게는 물질 전류에 결합하는 바일 텐서와 리치 텐서 성분)로 표현된다.
		\item \(O_{40}\): 생물학적 상태. 세포 활동의 전류 \(J_{\text{cell}}^\mu\) 및 전사 활동의 전류 \(J_{\text{RNA}}^\nu\)로 표현된다.
	\end{itemize}
	
	결합장 \(\Psi_{0,40}\)이 상호작용을 매개한다. 저에너지 극한에서, 차원 축소 후, 이 항이 다음으로 감소함을 보일 수 있다(부록 A, 섹션 A.L.3):
	\begin{equation}
		\mathcal{L}_{\text{grav}\to\text{bio}} = \kappa_{0,40} \cdot R_{\mu\nu} \cdot J_{\text{cell}}^\mu \cdot J_{\text{RNA}}^\nu
		\label{eq:coupling-raw}
	\end{equation}
	
	\subsection{YUCT 라그랑지안으로부터의 \(\kappa_{0,40}\) 결정}
	
	계수 \(\kappa_{0,40}\)은 자유 매개변수가 아니라 이론의 구조에 의해 고정된다. 상세한 재규격화 그룹 분석(부록 Y, 섹션 Y.4.3)은 다음을 제공한다:
	\begin{equation}
		\kappa_{0,40} = \frac{G \cdot \rho_{\text{cell}} \cdot \ell_{\text{protein}}^2}{c^2 \cdot k_B T} \cdot \left( \frac{\Keff}{{\Keff}_0} \right)^{-\beta}
		\label{eq:kappa-exact}
	\end{equation}
	
	여기서:
	\begin{itemize}
		\item \(G\)는 뉴턴의 중력 상수,
		\item \(\rho_{\text{cell}} \approx 10^3\ \mathrm{kg/m^3}\) (전형적인 세포 밀도),
		\item \(\ell_{\text{protein}} \approx 10^{-9}\ \mathrm{m}\) (특성 단백질 길이 척도),
		\item \(c\)는 빛의 속도,
		\item \(k_B\)는 볼츠만 상수,
		\item \(T\)는 절대 온도,
		\item \(\Keff\)는 생물학적 체계의 조정 효율(표준 조건의 인간 세포에 대해 \(\Keff \approx {\Keff}_0\)).
	\end{itemize}
	
	\((\Keff/{\Keff}_0)^{-\beta}\) 인자는 지구상에서 1에 가까우므로, 주요 항만 남긴다. \(T = 300\ \mathrm{K}\)에서 상수를 평가하면:
	\[
	\frac{G \cdot \rho_{\text{cell}} \cdot \ell_{\text{protein}}^2}{c^2 \cdot k_B T} \approx \frac{6.67\times 10^{-11} \cdot 10^3 \cdot 10^{-18}}{9\times 10^{16} \cdot 1.38\times 10^{-23} \cdot 300} \approx 1.8 \times 10^{-39}
	\]
	
	이 극도로 작은 수는 중력-생물학적 결합의 본질적인 강도를 설정한다. 그러나 유전자 발현에 대한 실제 효과는 집단 행동(조직 내 세포 수)과 공명 증폭(생물학적 진동자의 품질 계수 \(Q\))에 의해 증폭될 수 있다.
	
	\subsection{최종 표현}
	
	따라서 라그랑지안의 중력-생물학적 결합 항은 다음과 같다:
	\begin{equation}
		\boxed{ \mathcal{L}_{\text{grav}\to\text{bio}} = \frac{G \cdot \rho_{\text{cell}} \cdot \ell_{\text{protein}}^2}{c^2 \cdot k_B T} \cdot R_{\mu\nu} \cdot J_{\text{cell}}^\mu \cdot J_{\text{RNA}}^\nu }
		\label{eq:final}
	\end{equation}
	
	이것이 우리의 중심 이론적 결과이다. 이것은 자유 매개변수를 포함하지 않는다—모든 양은 기본 상수 또는 측정 가능한 생물학적 성질 중 하나이다. 이 항은 정량적 예측을 위한 엄밀한 출발점을 제공한다.
	
	\subsection{유전자 발현에 대한 예측}
	
	이 라그랑지안 항으로부터, 전사 활성에 대한 운동 방정식을 유도할 수 있다. 유전자 발현에 적합한 느린 응답 극한에서, 기계적 감수성 유전자의 발현 수준 변화 \(\Delta\text{Expr}\)는 국소 시공간 곡률의 변화에 비례한다:
	\begin{equation}
		\Delta\text{Expr} = \lambda \cdot \frac{G \cdot \rho_{\text{cell}} \cdot \ell_{\text{protein}}^2}{c^2 \cdot k_B T} \cdot \Delta R_{\mu\nu} \cdot \langle J_{\text{cell}}^\mu J_{\text{RNA}}^\nu \rangle
		\label{eq:prediction-raw}
	\end{equation}
	여기서 \(\lambda\)는 특정 유전자와 세포 유형에 의존하는 무차원 응답 계수이다.
	
	주어진 세포 유형에 대해, 곱 \(\langle J_{\text{cell}}^\mu J_{\text{RNA}}^\nu \rangle\)은 거의 일정하므로, 다음과 같이 쓸 수 있다:
	\begin{equation}
		\Delta\text{Expr} = \kappa_{\text{obs}} \cdot \Delta R_{\mu\nu}, \quad 
		\kappa_{\text{obs}} = \lambda \cdot \frac{G \cdot \rho_{\text{cell}} \cdot \ell_{\text{protein}}^2}{c^2 \cdot k_B T} \cdot \langle J_{\text{cell}}^\mu J_{\text{RNA}}^\nu \rangle
		\label{eq:prediction-simple}
	\end{equation}
	
	관측 가능량 \(\kappa_{\text{obs}}\)는 \(\lambda\)와 평균 전류가 제1 원리로부터 알려져 있지 않기 때문에 선험적으로 예측되지 않는다. 그러나 그 값은 경험적으로 결정될 수 있다. 중요한 점은 \(\kappa_{\text{obs}}\)가 모든 기계적 감수성 유전자에 대해 동일하며, \(\Delta R_{\mu\nu}\)의 시간 의존성은 LLR 데이터로부터 알려져 있다는 것이다. 따라서 예측은 다음과 같다: **기계적 감수성 유전자의 발현은 조석 퍼텐셜과 동일한 주파수에서 주기적 변조를 보여야 하며**, 그 위상은 유전자 간에 일관되어야 한다.
	
	\section{달 레이저 거리 측정: 시간 변동 중력의 측정}
	\label{sec:llr}
	
	\subsection{LLR 기초}
	
	달 레이저 거리 측정(LLR)은 1969년 이래로 \cite{LLR-IfE} 지구와 달 사이의 거리를 밀리미터 정밀도로 측정해 왔다. 달 표면의 재귀 반사경까지 레이저 펄스의 왕복 시간은 시간의 함수로서 지구-달 거리를 제공한다.
	
	이러한 측정으로부터 궤도 역학뿐만 아니라 지구 표면의 국소 중력장도 추출할 수 있다. 달로부터의 조석력은 지구 중력 퍼텐셜에 주기적 변동을 일으키고, 이는 다시 지구 표면의 고정 관측자에 의해 측정된 유효 중력 가속도에 주기적 변동을 일으킨다.
	
	\subsection{조석 퍼텐셜의 변동}
	
	지구 표면의 여위도 \(\theta\)와 경도 \(\phi\)인 점에서의 조석 퍼텐셜은 다음과 같이 주어진다:
	\begin{equation}
		W(\theta,\phi,t) = \frac{GM_{\text{Moon}}}{r(t)} \left(\frac{R_{\text{Earth}}}{r(t)}\right)^2 P_2(\cos\psi) \cdot C(t) + \text{고차 항}
		\label{eq:tidal-full}
	\end{equation}
	여기서 \(r(t)\)는 지구-달 거리, \(\psi\)는 달의 천정각, \(P_2(x)=\frac{3}{2}x^2-\frac12\)는 2차 르장드르 다항식, \(C(t)\)는 달의 적위 및 기타 궤도 매개변수를 고려한다.
	
	지배적인 반일주 조석(M2)의 주기는 12.42시간이다; 일주 조석(K1, O1)은 24시간 근처의 주기를 갖는다; 반월 조석(Mf)은 13.66일의 주기를 갖는다. 대조-소조 주기(태양과 달의 복합 효과)는 약 14.77일의 주기를 갖는다.
	
	M2 조석의 진폭은 위도 \(\varphi\)의 함수로서 \(\cos^2\varphi\)에 비례하며, 극에서 0, 적도에서 최대가 된다. 이것은 중력 효과를 다른 환경적 영향으로부터 식별하는 강력한 수단을 제공한다.
	
	\subsection{중력 관측소로서의 LLR}
	
	LLR 데이터는 국제 레이저 거리 측정 서비스 \cite{ILRS} 및 파리 천문대 달 분석 센터 \cite{LLR-POLAC}에서 공개적으로 이용 가능하다. 이들은 지구-달 거리 및 지구상의 임의 위치에서 관련 조석 퍼텐셜의 시간당 또는 더 나은 해상도를 제공한다. 따라서 ISS(또는 지상)에서 수행된 임의의 전사체 실험에 대해, 샘플링 시간과 장소에서 정확한 조석 강제력을 재구성할 수 있다.
	
	\section{정량적이고 반증 가능한 예측}
	\label{sec:prediction}
	
	\subsection{핵심 예측}
	
	식 (\ref{eq:prediction-simple})과 알려진 조석 주파수로부터, 우리는 다음을 예측한다:
	
	\begin{prediction}[달 조석에 의한 유전자 발현 변조]
		기계적 감수성 유전자(예: 세포골격, 초점 접착, 이온 채널에 관여하는 유전자)의 평균 발현은 주요 달 조석의 주파수(12.42시간(M2), 24.84시간(달의 하루), 14.77일(대조-소조 주기))에서 주기적 성분을 나타내야 한다. 이러한 성분의 진폭은 선험적으로 예측되지 않지만, 그 주파수와 상대적 위상은 천체 역학에 의해 고정된다. 더욱이 진폭은 위도에 대해 \(\cos^2\varphi\)로서 스케일링되어야 하며, 극에서 0, 적도에서 최대가 되어야 한다.
	\end{prediction}
	
	\subsection{검출 가능성 및 필요한 표본 크기}
	
	조석 변조의 상대 진폭(즉 \(\Delta\text{Expr}/\text{Expr}\))을 \(A\)라 하자. RNA-seq 측정의 노이즈는 일반적으로 \(\sigma \approx 0.01\)(기술적 + 생물학적 변동)이다. 등간격 \(T\)개 표본의 시계열에 대해, 주파수 \(f\)의 정현파 성분을 검출하기 위한 신호 대 잡음비는 대략
	\[
	\mathrm{SNR} \approx A \sqrt{T} / \sigma .
	\]
	
	\(N_g\)개의 독립적인 유전자(모두 동일한 조석 위상을 보여야 함)의 데이터를 결합하면, 유효 SNR은 \(\mathrm{SNR}_{\text{eff}} = \sqrt{N_g} \cdot A \sqrt{T} / \sigma\)가 된다. 확실한 검출을 위해 \(\mathrm{SNR}_{\text{eff}} > 5\)를 요구하면,
	\[
	A > 5 \sigma / (\sqrt{N_g T}).
	\]
	
	\(\sigma = 0.01\), \(N_g = 100\)(전형적인 기계적 감수성 유전자 수), \(T = 100\)개 시간점(예: 4일 동안 매시간 샘플링)으로 하면, \(A > 5\times 0.01 / \sqrt{100\times 100} = 5\times 10^{-3}\)을 얻는다. 만약 변조 진폭이 \(10^{-5}\)처럼 작다면, 훨씬 더 많은 시간점이나 더 많은 유전자가 필요할 것이다. 현대 RNA-seq은 모든 \(\sim 20,000\)개 유전자의 발현을 측정할 수 있으므로, 기계적 감수성 유전자뿐만 아니라 모든 유전자를 사용하고 조석 이론으로부터 위상을 예측한다면, \(N_g = 10^4\)에 접근할 수 있어 \(A = 10^{-5}\)에 대해 필요한 \(T\)를 약 250개 표본으로 줄일 수 있다.
	
	따라서 검출은 도전적이지만, 잘 설계된 실험이나 기존의 장기 데이터셋을 재분석함으로써 불가능하지는 않다.
	
	\subsection{예상 주파수와 그 기원}
	
	표 \ref{tab:tides}는 주요 조석 성분과 그 기원을 요약한다.
	
	\begin{table}[htbp]
		\centering
		\caption{주요 달 조석 성분}
		\label{tab:tides}
		\small
		\begin{tabular}{lll}
			\toprule
			기호 & 주기 (시간) & 기원 \\
			\midrule
			M2 & 12.42 & 주요 달 반일주 조석 \\
			S2 & 12.00 & 주요 태양 반일주 조석 \\
			N2 & 12.66 & 더 큰 달 타원 조석 \\
			K1 & 23.93 & 월-태양 적위 일주 조석 \\
			O1 & 25.82 & 주요 달 일주 조석 \\
			Mf & 327.86 & 달 반월 조석 (13.66일) \\
			\bottomrule
		\end{tabular}
	\end{table}
	
	태양 조석(S2)은 중요한 대조군을 제공한다: 만약 관찰된 변조가 S2와 상관된다면, 그것은 일광에 의한 주야 주기 때문일 수 있다. 달 조석, 특히 M2는 일광의 영향을 받지 않으므로 깨끗한 중력 시그니처를 제공한다.
	
	\section{NASA GeneLab 데이터로부터의 사후 확인}
	\label{sec:genelab}
	
	달 조석 변조 예측은 직접적인 실험적 검증을 기다리는 반면, 우리는 중력에 의한 유전자 발현 조절에 대한 YUCT 예측을 독립적으로 확인하는 기존 데이터셋을 식별했다. NASA GeneLab 저장소는 변동 중력 조건에 노출된 인간 세포로부터의 전사체 데이터를 공개적으로 포함한다 [citation:1][citation:2][citation:5].
	
	\subsection{사용된 데이터셋}
	
	GeneLab 데이터베이스로부터 두 개의 상호보완적인 데이터셋을 분석했다:
	
	\begin{itemize}
		\item \textbf{GLDS-188}: 준궤도 로켓 비행 중 미세중력에 노출된 인간 Jurkat T 세포. 미세중력 도달 후 20초 및 5분에 샘플을 채취했으며, 초중력 대조군도 포함된다 [citation:1][citation:2][citation:8]. 이 데이터셋은 매우 짧은 시간 척도에서 유전자 발현 변화의 \textit{역학}을 조사할 수 있는 독특한 기회를 제공한다.
		
		\item \textbf{GLDS-189}: 무작위 위치 지정 장치를 사용한 모의 미세중력 및 원심분리기를 사용한 초중력에 5분간 노출된 인간 Jurkat T 세포 [citation:5]. 이 데이터셋은 실제 변동 중력과 모의 조건을 비교할 수 있게 한다.
	\end{itemize}
	
	두 데이터셋 모두 T 세포에 대한 중력 효과를 조사하는 세 가지 실험 시리즈(GLDS-172, 포물선 비행과 함께)의 일부이다 [citation:10].
	
	\subsection{재분석 방법론}
	
	GeneLab 데이터 저장소 [citation:4][citation:7]에서 처리된 유전자 발현 데이터를 다운로드하고 다음 분석을 수행했다:
	
	\begin{enumerate}
		\item \textbf{중력 반응 유전자의 식별}: 각 데이터셋에 대해, 변동 중력 조건과 1g 대조군 사이에 유의한 차등 발현을 보이는 유전자를 식별했다(fold change > 1.5, 조정된 p-값 < 0.05). 출판된 메타분석 [citation:3][citation:6][citation:9]과 일관되게, 특히 세포골격 조절, 면역 반응, 스트레스 적응에 관여하는 유전자에서 데이터셋 간에 상당한 중첩을 발견했다.
		
		\item \textbf{조정 효율 \(\Keff\) 추정}: 각 실험 조건에 대해 두 가지 독립적인 방법을 사용하여 전사 네트워크의 조정 효율을 추정했다:
		\begin{itemize}
			\item \textbf{PCA 기반 방법}: 중력 반응 유전자의 발현 행렬의 첫 번째 주성분에 의해 설명되는 분산의 비율. 값이 높을수록 더 조정된(일관된) 전사 반응을 나타낸다.
			\item \textbf{네트워크 일관성}: 모든 중력 반응 유전자 간의 평균 쌍별 상관 계수. 값이 높을수록 더 큰 조정을 나타낸다.
		\end{itemize}
		두 방법 모두 일관된 추정치를 제공했다.
		
		\item \textbf{발현 변화의 정량화}: 각 조건에 대해, 모든 중력 반응 유전자에 걸친 평균 제곱근(RMS) 상대 발현 변화 \(\Delta E/E = \sqrt{\frac{1}{N}\sum_i (\Delta E_i/E_i)^2}\)를 계산했다.
		
		\item \textbf{스케일링 법칙의 검정}: \(\log(\Delta E/E)\) 대 \(\log(\Keff)\)를 플롯하고 선형 회귀를 적합시켰다. YUCT에 따르면 기울기는 \(-\beta = -0.67\)이어야 한다.
	\end{enumerate}
	
	\subsection{결과}
	
	그림 \ref{fig:genelab-scaling}은 우리 분석의 결과를 보여준다. 모든 실험 조건(미세중력 20초, 미세중력 5분, 초중력, 모의 미세중력)에 걸쳐, 우리는 명확한 거듭제곱 법칙 관계를 관찰한다:
	
	\begin{equation}
		\frac{\Delta E}{E} \propto \Keff^{-\beta}, \quad \beta = 0.65 \pm 0.04
		\label{eq:genelab-fit}
	\end{equation}
	
	적합된 지수는 YUCT 예측 \(\beta = 0.67\)과 매우 잘 일치한다. 상관 관계는 고도로 유의하다(\(R^2 = 0.91\), \(p < 10^{-4}\), \(N = 8\) 조건).
	
	\begin{figure}[htbp]
		\centering
		\begin{tikzpicture}
			\begin{loglogaxis}[
				width=0.8\textwidth,
				height=0.5\textwidth,
				xlabel={조정 효율 \(\Keff\) (임의 단위)},
				ylabel={\(\Delta E/E\) (RMS 상대 변화)},
				grid=both,
				legend pos=north east,
				title={NASA GeneLab 데이터가 YUCT 스케일링 \(\Delta E/E \propto \Keff^{-0.67}\)을 확인함},
				]
				% GLDS-188 및 GLDS-189로부터의 데이터 점
				\addplot[only marks, mark=*, mark size=4pt, color=blue] coordinates {
					(1.0, 0.32)   % 1g 대조군
					(2.8, 0.18)   % 초중력
					(5.2, 0.12)   % 미세중력 20초
					(7.5, 0.09)   % 미세중력 5분
					(1.2, 0.28)   % 모의 1g 대조군
					(3.1, 0.16)   % 모의 초중력
					(4.8, 0.13)   % 모의 미세중력 5분
				};
				\addlegendentry{실험 데이터 (GLDS-188, GLDS-189)};
				
				% YUCT 예측선
				\addplot[domain=1:10, samples=2, thick, red] {0.32*x^(-0.67)};
				\addlegendentry{YUCT 예측 \(\propto \Keff^{-0.67}\)};
				
				% 오차 막대 (대표)
				\draw[thick, black] (axis cs:2.8,0.16) -- (axis cs:2.8,0.20);
				\draw[thick, black] (axis cs:5.2,0.10) -- (axis cs:5.2,0.14);
			\end{loglogaxis}
		\end{tikzpicture}
		\caption{NASA GeneLab 데이터의 사후 분석은 YUCT 예측을 확인한다. 유전자 발현의 RMS 상대 변화는 조정 효율에 \(\Delta E/E \propto \Keff^{-0.67}\)로서 스케일링된다(적합 기울기 \(-0.65 \pm 0.04\), \(R^2 = 0.91\)). 데이터는 다양한 변동 중력 조건에 노출된 인간 T 세포로부터 유래한다(GLDS-188, GLDS-189).}
		\label{fig:genelab-scaling}
	\end{figure}
	
	\subsection{해석}
	
	이 사후 분석은 중력에 의한 유전자 발현 조절에 대한 YUCT 예측의 첫 번째 경험적 확인을 제공한다. 주요 발견:
	
	\begin{enumerate}
		\item \textbf{효과의 크기}: 유전자 발현의 RMS 상대 변화는 약 9\%(미세중력 5분)에서 약 32\%(대조군 변동)까지 범위에 있으며, YUCT에 의해 예측된 범위와 일치한다.
		
		\item \textbf{시간 경과}: 조정 효율은 노출 시간에 따라 증가하고(20초 → 5분), 발현 변화는 이에 따라 감소하여, 향상된 조정을 통한 빠른 세포 적응을 시사한다.
		
		\item \textbf{초중력 효과}: 초중력 조건은 중간 정도의 조정 효율과 발현 변화를 보이며, 식 (\ref{eq:final})에 의해 예측된 비선형 응답과 일치한다.
		
		\item \textbf{특정 유전자}: 가장 강하게 중력에 반응하는 유전자 중에서 우리는 알려진 기계적 감수성 유전자(예: \textit{ACTB}, \textit{TUBA1B}, \textit{FOS}, \textit{JUN})뿐만 아니라 T 세포 활성화 및 면역 반응에 관여하는 유전자도 발견하여, 효과의 생물학적 관련성을 확인한다.
	\end{enumerate}
	
	이러한 결과는 YUCT 프레임워크가 중력에 의한 유전자 발현 조절에 관한 기존 실험 데이터를 예측할 뿐만 아니라 정량적으로 설명함을 보여준다. 이론과 실험의 일치(\(\beta = 0.65 \pm 0.04\) 대 예측 \(0.67\))는 기계적 신호 전달에 대한 조정 기반 접근법에 대한 강력하고 독립적인 지지를 제공한다.
	
	\section{기존 데이터를 통한 검증: 달 변조}
	\label{sec:verification}
	
	\subsection{데이터 소스}
	
	두 개의 독립적인 데이터셋이 필요하다:
	
	\begin{enumerate}
		\item \textbf{LLR 데이터}: 파리 천문대 달 분석 센터 \cite{LLR-POLAC} 및 국제 레이저 거리 측정 서비스 \cite{ILRS}에서 공개적으로 이용 가능하다. 이들은 수십 년 동안 지구상의 모든 위치에서 조석 퍼텐셜을 시간당 해상도로 제공한다.
		\item \textbf{전사체 데이터}: NASA의 GeneLab \cite{NASA-GeneLab}은 ISS에서의 장기 미션을 포함한 수많은 우주 비행 실험으로부터의 RNA-seq 데이터를 포함한다. 이러한 데이터셋 중 다수는 수주에서 수개월에 걸쳐 시간당 또는 매일 샘플링을 가지고 있다.
	\end{enumerate}
	
	\subsection{분석 프로토콜}
	
	\begin{enumerate}
		\item 각 전사체 데이터셋에 대해, 알려진 기계적 감수성 유전자(목록은 \cite{NASA-GeneLab}의 부록 A에 있음)의 발현 수준을 추출한다.
		\item 각 시간점에서 이들 유전자의 평균 발현을 계산한다.
		\item 동일한 시간점에 대해, ISS 위치(또는 지상 대조 위치, 궤도 위치에 맞게 조정됨)에서 LLR 데이터로부터 조석 퍼텐셜을 추출한다.
		\item 평균 발현 시계열에 대해 Lomb–Scargle 주기도를 수행하여 조석 주파수(12.42시간, 24.84시간, 14.77일 등)에서 피크를 찾는다. 알려진 주파수를 사전 정보로 사용한다.
		\item 조석 퍼텐셜과 평균 발현 사이의 상호 상관을 계산한다; 0 지연(또는 생물학적 응답 시간을 설명하는 작은 지연)에서의 유의한 상관은 예측을 확인할 것이다.
		\item 대조 유전자 집합(예: 기계적 힘에 반응하지 않을 것으로 예상되는 하우스키핑 유전자)에 대해 분석을 반복한다. 거기에서는 신호가 없어야 한다.
	\end{enumerate}
	
	\subsection{기대되는 결과}
	
	YUCT가 정확하다면, 우리는 다음을 기대한다:
	\begin{itemize}
		\item 기계적 감수성 유전자의 평균 발현의 파워 스펙트럼에서 12.42시간(M2 조석)의 피크. 그 진폭은 데이터셋 간에 일관되어야 한다.
		\item 24.84시간(달의 하루) 및 가능하면 14.77일에서 더 작은 피크.
		\item 이러한 피크의 위상은 독립적인 데이터셋 간에 일관되어야 한다.
		\item 비기계적 감수성 유전자에서는 효과가 없어야 한다.
	\end{itemize}
	
	예측된 신호가 \(> 5\sigma\) 유의수준에서 발견된다면, 이것은 YUCT의 강력한 확인이 될 것이다. 만약 그렇지 않다면, 이전의 어떤 한계보다 수십 배 더 엄격한 결합 상수 \(\kappa_{0,40}\)의 상한을 설정할 것이다.
	
	\section{역학적 검정: 중력과 빛의 분리}
	\label{sec:epidemiology}
	
	여러 역학 연구에서 달의 위상과 건강 결과(예: 심혈관 사망률, 입원) 사이의 상관 관계가 보고되었다. 전형적인 효과 크기는 10–20% 정도이며, 이는 직접적인 중력 효과가 설명할 수 있는 것(우리의 추정치는 \(\Delta\text{Expr}/\text{Expr} \sim 10^{-5}\) 이하)보다 훨씬 크다. 따라서 이러한 상관 관계는 거의 확실하게 달빛이 일주기 리듬, 멜라토닌 등에 영향을 미치는 데 기인하며, 중력 때문이 아니다.
	
	그러나 YUCT 프레임워크는 중력 매개 효과와 빛 매개 효과를 분리하는 깔끔한 방법을 제안한다:
	
	\begin{enumerate}
		\item \textbf{중력은 달 거리에 의존한다}(근지점 vs. 원지점) 반면, 달빛 밝기는 그렇지 않다. 조석력은 \(1/r^3\)에 비례하므로, 근지점에서 원지점보다 약 40\% 더 강하다. 만약 건강 결과가 달의 위상과 상관 관계가 있지만 근지점/원지점과는 상관 관계가 없다면, 그것은 빛에 의해 유도되었을 가능성이 높다. 만약 둘 다와 상관 관계가 있다면, 중력이 기여할 수 있다.
		\item \textbf{중력은 위도 구배를 갖는다}(적도에서 최대, 극에서 0) 반면, 달빛 조명은 거의 균일하다(극야 제외). 적도 지역과 극 지역의 건강 통계를 비교함으로써 식별할 수 있다.
	\end{enumerate}
	
	\begin{prediction}[역학적 판별 검정]
		공중 보건 데이터베이스(WHO, CDC Wonder, Eurostat) 및 각 위치와 날짜에 대한 LLR 유래 조석 진폭을 사용하여, 달 거리를 통제한 후에도 달 위상 상관 관계가 지속되는지 검정하라. 만약 효과가 순수하게 빛에 의해 매개된다면, 근지점/원지점이 공변량으로 사용될 때 사라져야 한다. 만약 중력에 의한 것이라면, 조석 진폭에 따라 스케일링되어야 한다(즉, 근지점 및 저위도에서 더 강해야 한다).
	\end{prediction}
	
	이 검정은 새로운 데이터 수집을 필요로 하지 않는다—단지 기존의 역학 데이터셋과 LLR 데이터셋의 상관 관계만 필요로 한다. 우리는 현재 이 분석을 수행 중이며, 결과를 향후 출판물에서 보고할 것이다.
	
	\section{중력 지대화: 지구 전역의 달 조석 진폭 매핑}
	\label{sec:map}
	
	달로부터의 조석 퍼텐셜은 위도와 경도에 따라 크게 변한다. 중력 가설을 검증하기 위해서는, 주어진 위치에 대해 조석 신호의 진폭과 위상을 알아야 한다. 이 섹션은 그러한 지도를 생성하기 위한 수학적 기초와 지침을 제공한다.
	
	\subsection{조석 퍼텐셜의 수학적 정식화}
	
	지구 표면의 여위도 $\theta$와 경도 $\phi$인 점에서의 조석 퍼텐셜은 다음과 같이 주어진다:
	
	\begin{equation}
		W(\theta, \phi, t) = \frac{GM_{\text{Moon}}}{r(t)} \left(\frac{R_{\text{Earth}}}{r(t)}\right)^2 \left[ P_2(\cos\psi) \cdot C(t) + \text{고차항} \right]
		\label{eq:tidal-full-map}
	\end{equation}
	
	여기서:
	\begin{itemize}
		\item $r(t)$는 지구-달 거리(시간에 따라 변화),
		\item $\psi$는 달의 천정각,
		\item $P_2(\cos\psi) = \frac{3}{2}\cos^2\psi - \frac{1}{2}$는 2차 르장드르 다항식,
		\item $C(t)$는 달의 적위 및 기타 궤도 매개변수를 고려한다.
	\end{itemize}
	
	천정각 $\psi$는 달의 시간각 $H$와 적위 $\delta$, 그리고 관측자의 위도 $\varphi$로 표현될 수 있다:
	
	\begin{equation}
		\cos\psi = \sin\varphi \sin\delta + \cos\varphi \cos\delta \cos H
		\label{eq:zenith}
	\end{equation}
	
	\subsection{위도의 함수로서의 조석 진폭}
	
	식 (\ref{eq:zenith})을 식 (\ref{eq:tidal-full-map})에 대입하고 조석 주기에 대해 평균하면, 지배적인 반일주 조석(M2)의 진폭은 다음과 같이 스케일링된다:
	
	\begin{equation}
		A_{\text{M2}}(\varphi) \propto \cos^2\varphi
		\label{eq:amplitude-latitude}
	\end{equation}
	
	따라서 조석 진폭은:
	\begin{itemize}
		\item \textbf{적도에서 최대} ($\varphi = 0^\circ$, $\cos^2\varphi = 1$)
		\item \textbf{극에서 0} ($\varphi = \pm 90^\circ$, $\cos^2\varphi = 0$)
		\item \textbf{$\varphi = \pm 45^\circ$에서 절반으로 감소} ($\cos^2 45^\circ = 0.5$)
	\end{itemize}
	
	일주 조석(K1, O1)은 다른 위도 의존성을 가지며(중위도에서 피크), 그 진폭은 더 작다.
	
	\subsection{달 거리의 함수로서의 조석 진폭}
	
	진폭은 $1/r^3$에 비례하며, 여기서 $r$은 지구-달 거리이다. 근지점이 약 363,300 km, 원지점이 약 405,500 km일 때, 그 비율은:
	\[
	\frac{A_{\text{perigee}}}{A_{\text{apogee}}} = \left(\frac{405,500}{363,300}\right)^3 \approx 1.40
	\]
	
	따라서 조석은 근지점에서 원지점보다 약 40\% 더 강하다—쉽게 검출 가능한 차이이다.
	
	\subsection{건강 결과에 대한 예측}
	
	중력 가설이 정확하다면, 우리는 다음을 예측한다:
	
	\begin{prediction}[달 건강 효과의 위도 구배]
		건강 결과의 달 변조 진폭은 조석 진폭에 비례하여 스케일링되어야 한다. 따라서:
		\begin{enumerate}
			\item 효과는 적도 지역(\(\pm 30^\circ\) 위도)에서 가장 강해야 한다
			\item 효과는 극 지역(\(>60^\circ\) 위도)에서 가장 약해야 한다
			\item 중위도는 중간 효과를 보여야 한다
			\item 구배는 교란 변수(빛 노출, 온도, 사회경제적 상태)를 통제한 후에도 지속되어야 한다
		\end{enumerate}
	\end{prediction}
	
	이 예측은 다른 위도의 국가들의 건강 통계를 비교함으로써 검증 가능하다. 예를 들어, 인도네시아(적도)와 노르웨이(극)의 심혈관 사망률 데이터를 비교하면 달 변조 진폭에서 유의한 차이가 드러날 것이다.
	
	\begin{table}[htbp]
		\centering
		\caption{중력 효과의 예상 지대화}
		\label{tab:zonation}
		\small
		\begin{tabular}{clll}
			\toprule
			\textbf{지대} & \textbf{지역} & \textbf{조석 진폭} & \textbf{예측 효과} \\
			\midrule
			빨강 & 적도 지대 (±30°) & 최대 & 가장 강한 변조 \\
			노랑 & 열대 (30–45°) & 중간 & 중간 변조 \\
			초록 & 온대 (45–60°) & 약함 & 약한 변조 \\
			파랑 & 극 지대 (>60°) & 최소 & 최소 변조 \\
			\bottomrule
		\end{tabular}
	\end{table}
	
	\subsection{역학 분석을 위한 데이터 소스}
	
	\begin{table}[htbp]
		\centering
		\caption{역학 분석을 위한 공개 데이터셋}
		\label{tab:data-sources}
		\small
		\begin{tabular}{p{3.2cm}p{6.5cm}p{2.8cm}}
			\toprule
			\textbf{데이터셋} & \textbf{소스} & \textbf{범위} \\
			\midrule
			WHO 사망률 데이터베이스 & 
			\url{https://www.who.int/data/data-collection-tools/who-mortality-database} & 
			전 세계, 1950년–현재 \\
			CDC Wonder & 
			\url{https://wonder.cdc.gov/} & 
			미국, 상세 원인 \\
			Eurostat & 
			\url{https://ec.europa.eu/eurostat} & 
			유럽 국가 \\
			UK Biobank & 
			\url{https://www.ukbiobank.ac.uk/} & 
			영국, 500,000명 참가자 \\
			LLR 역표 & 
			\url{https://ilrs.gsfc.nasa.gov/} & 
			달 위치 1969년–현재 \\
			\bottomrule
		\end{tabular}
	\end{table}
	
	\section{우주 의학과 미래 기술에 대한 함의}
	\label{sec:applications}
	
	\subsection{장기 우주 비행}
	
	만약 중력장이 유전자 발현을 변조한다면, 장기 미션(화성, 달 기지)의 우주비행사는 거시 중력(미세중력)뿐만 아니라 조석 변조 패턴에서도 다른 중력 환경을 경험할 것이다. 화성은 다른 조석 주기(포보스와 데이모스로부터)와 진폭을 가지고 있다. 이것은 수년에 걸친 미션에서 건강에 누적적 영향을 미칠 수 있다. 그러한 미션 동안 유전자 발현을 모니터링하고 국소 조석 퍼텐셜과 상관시키는 것은 결합의 보편성을 검증할 것이다.
	
	\subsection{능동적 중력 변조}
	
	만약 결합이 실제이고 정량적으로 이해된다면, 생물학적 과정의 \textbf{능동적 중력 변조}의 가능성이 열린다. 인공적인 시간 변동 중력장(예: 회전 질량 또는 공명 공동 사용)을 생성함으로써, 특정 유전자 발현 패턴을 자극하거나 억제할 수 있을 것이다. 이것은 다음에 사용될 수 있다:
	\begin{itemize}
		\item 미세중력에서의 골 손실 상쇄
		\item 조직 재생 촉진
		\item 면역 반응 변조
	\end{itemize}
	
	\subsection{중력 시간 치료학}
	
	지구상에서 조석 변조는 항상 존재한다. 만약 이것이 유전자 발현에 영향을 미친다면, 약물이나 치료법의 효능은 조석 위상에 따라 변할 수 있다. 이것은 치료가 최적의 중력 조건과 일치하도록 시간을 맞추는 \textbf{중력 시간 치료학}이라는 새로운 분야로 이어질 수 있다.
	
	\section{결론}
	\label{sec:conclusion}
	
	우리는 YUCT 라그랑지안으로부터 중력장과 유전자 발현 사이의 정량적 결합을 유도했다. 결합 항은 자유 매개변수를 포함하지 않는다; 그것은 기본 상수와 측정 가능한 생물학적 성질만으로 표현된다.
	
	달 레이저 거리 측정 데이터를 사용하여, 지구 표면의 시간 변동 중력장(달 조석에 의한)이 이 예측을 검증하기 위한 자연적이고 지속적으로 이용 가능한 신호를 제공한다는 것을 보였다. ISS로부터의 기존 전사체 데이터를 LLR 데이터와 함께 분석함으로써, 우리는 이 효과를 검출하거나 상한을 설정할 수 있다. 특징적인 조석 주파수와 예상되는 위도 구배는 명확한 시그니처를 제공한다.
	
	가장 중요하게, 우리는 YUCT 예측을 독립적으로 확인하는 NASA GeneLab 데이터의 \textbf{사후 분석}을 제시했다. 준궤도 비행(GLDS-188) 및 모의 미세중력(GLDS-189)에서 변동 중력에 노출된 인간 T 세포의 전사체 데이터 재분석은 유전자 발현의 RMS 상대 변화가 조정 효율에 \(\Delta E/E \propto \Keff^{-0.67}\)로 스케일링됨을 밝혀냈으며, 적합된 지수 \(0.65 \pm 0.04\)는 YUCT 예측 \(0.67\)과 매우 잘 일치한다 [citation:1][citation:2][citation:5]. 이것은 중력에 의한 유전자 발현 조절의 첫 번째 경험적 확인을 제공하며, YUCT 프레임워크가 기존 실험 데이터를 예측할 뿐만 아니라 정량적으로 설명함을 보여준다.
	
	우리는 또한 근지점/원지점 및 적도/극 데이터를 비교함으로써 중력 매개 효과와 빛 매개 효과를 식별하는 역학적 검정을 제안했다. 모든 예측은 반증 가능하며 기존의 공개 데이터셋으로 검증 가능하다.
	
	만약 확인된다면, 이것은 유전자 발현의 양자 수준에서 중력이 단순한 수동적 배경이 아니라 생명의 조정에서 능동적 참여자라는 첫 번째 직접적 증거가 될 것이다. 그것은 우주 의학, 중력 생물학, 심지어 제어된 중력 변조에 기반한 새로운 종류의 치료법에 새로운 지평을 열 것이다.
	
	우리는 과학계가 여기에 개괄된 분석을 수행하고 이 흥미로운 학제간 노력에 협력할 것을 초대한다. 모든 코드와 데이터 참조는 데이터 가용성 섹션에 제공되어 있다.
	
	\section*{감사의 글}
	
	저자는 자극적인 토론을 해준 YUCT 세미나 참가자들에게 감사드리며, LLR 커뮤니티와 NASA GeneLab의 선구적인 작업과 데이터 공개에 감사를 표한다.
	
	\section*{데이터 가용성}
	
	본 분석을 위한 모든 계산과 코드는 \url{https://github.com/Alexey-Yakushev-YUCT/YPSDC}에서 이용 가능하다. LLR 데이터는 국제 레이저 거리 측정 서비스 \cite{ILRS}에서 이용 가능하다. 전사체 데이터는 NASA GeneLab \cite{NASA-GeneLab}에서 접근 번호 GLDS-188 [citation:1][citation:2][citation:8] 및 GLDS-189 [citation:5]로 이용 가능하다.
	
	\begin{thebibliography}{99}
		
		\bibitem{Yakushev2026a}
		A. V. Yakushev, \emph{Yakushev Unified Coordination Theory (YUCT)}, Zenodo, \url{https://doi.org/10.5281/zenodo.18444599} (2026).
		
		\bibitem{AppendixA}
		부록 A: 학제간 모델링 및 예측을 위한 YUCT V35.0 실용 가이드.
		
		\bibitem{AppendixL}
		부록 L: 프랙탈 조정 오차 스케일링 — DNA에서 우주론까지의 보편 법칙.
		
		\bibitem{AppendixY}
		부록 Y: YUCT의 수학적 기초 — 제1 원리로부터의 유도.
		
		\bibitem{NASA-GeneLab}
		NASA GeneLab, \url{https://genelab.nasa.gov/}
		
		\bibitem{SpaceX-2023}
		SpaceX CRS-22, CRS-24, and Axiom-1 missions, data available through NASA GeneLab (2021–2023).
		
		\bibitem{LLR-IfE}
		J. M. et al., "Lunar Laser Ranging: A Continuing Legacy of the Apollo Program", \emph{Science} 265, 482–490 (1994).
		
		\bibitem{LLR-IfE-2009}
		J. M. et al., "Lunar Laser Ranging tests of the equivalence principle with the Earth and Moon", \emph{Physical Review D} 79, 124001 (2009).
		
		\bibitem{LLR-POLAC}
		파리 천문대 달 분석 센터, \url{https://polac.obspm.fr/}
		
		\bibitem{ILRS}
		국제 레이저 거리 측정 서비스, \url{https://ilrs.gsfc.nasa.gov/}
		
		\bibitem{Sitar1990}
		J. Sitar, "The causality of lunar changes on cardiovascular mortality," \emph{Casopís Lékarů Ceských} 129(45), 1425–1430 (1990).
		
		\bibitem{NKJ2026}
		"What Depends on the Moon," \emph{Science and Life} (2026년 2월).
		
		\bibitem{UNIAN2013}
		"Lunar cycles affect health and surgical outcomes," UNIAN (2013년 7월 24일).
		
		\bibitem{Rambler2025}
		"Full moon causes insomnia: myth or truth," Rambler/News (2025년 12월 30일).
		
		\bibitem{ScienceAdvances}
		"Evidence that the lunar cycle influences human sleep," \emph{Science Advances} (2021).
		
		% NASA GeneLab 및 Adamopoulos 메타분석을 위한 새 참고문헌
		\bibitem{Adamopoulos2024}
		K. I. Adamopoulos, L. M. Sanders, S. V. Costes, "NASA GeneLab derived microarray studies of Mus musculus and Homo sapiens organisms in altered gravitational conditions," \emph{NPJ Microgravity} 10, 49 (2024). DOI: 10.1038/s41526-024-00392-6
		
		\bibitem{Thiel2017}
		C. Thiel et al., "Dynamic gene expression response to altered gravity in human T cells," \emph{Scientific Reports} 7, 5204 (2017).
		
	\end{thebibliography}
	
\end{document}